\documentclass[11pt,a4paper]{article}
\usepackage[utf8]{inputenc}
\usepackage{amsmath,amssymb,amsthm}
\usepackage{geometry}
\geometry{margin=2.5cm}
\usepackage{listings}
\usepackage{xcolor}
\usepackage{hyperref}

\lstset{
  language=C++,
  basicstyle=\ttfamily\small,
  keywordstyle=\color{blue}\bfseries,
  commentstyle=\color{green!60!black},
  stringstyle=\color{red!70!black},
  numbers=left,
  numberstyle=\tiny\color{gray},
  numbersep=8pt,
  frame=single,
  breaklines=true,
  tabsize=4,
  showstringspaces=false
}

\title{IOI 1996 -- Network of Schools}
\author{Solution}
\date{}

\begin{document}
\maketitle

\section{Problem Statement}

A network of $n$ schools is connected by one-way links. Software distributed to a school
propagates transitively along outgoing links.

\begin{itemize}
  \item \textbf{Task A:} Find the minimum number of schools that must receive the software
    initially so that all schools eventually receive it.
  \item \textbf{Task B:} Find the minimum number of one-way links to add so that
    distributing software to \emph{any single} school causes all schools to receive it
    (i.e., make the graph strongly connected).
\end{itemize}

\noindent\textbf{Constraints:} $2 \le n \le 100$.

\section{Solution Approach}

\subsection{Step 1: Strongly Connected Components}

Use Kosaraju's algorithm (or Tarjan's) to compute the SCCs. Condense the graph into a
DAG where each node represents an SCC.

\subsection{Step 2: Analyze the Condensed DAG}

Let $S$ denote the number of SCCs.

\begin{itemize}
  \item \textbf{Task A:} The answer is the number of \emph{source nodes}
    (in-degree 0) in the condensed DAG. These SCCs cannot be reached from any
    other SCC, so each must receive software directly:
    \[
      \text{Answer A} = |\{v : \deg^-(v) = 0\}|.
    \]

  \item \textbf{Task B:} To make the condensed DAG strongly connected, we must
    eliminate all sources and sinks. The minimum number of edges required is:
    \[
      \text{Answer B} = \max(\text{sources}, \text{sinks}),
    \]
    where sources and sinks refer to nodes with in-degree 0 and out-degree 0
    in the condensed DAG, respectively.
\end{itemize}

\begin{proof}[Justification for Task B]
Each added edge can eliminate at most one source and one sink. Therefore
$\max(\text{sources}, \text{sinks})$ is a lower bound. A matching construction
(pairing sources with sinks and connecting them cyclically) achieves this bound.
\end{proof}

\noindent\textbf{Special case:} If $S = 1$ (the graph is already strongly connected),
then Answer A = 1 and Answer B = 0.

\section{C++ Solution}

\begin{lstlisting}
#include <cstdio>
#include <cstring>
#include <vector>
#include <algorithm>
using namespace std;

const int MAXN = 105;
int n;
vector<int> adj[MAXN], radj[MAXN];
int comp[MAXN];
bool visited[MAXN];
int order_arr[MAXN], order_cnt;

// Kosaraju's Algorithm -- Pass 1: forward DFS, record finish order
void dfs1(int u) {
    visited[u] = true;
    for (int v : adj[u])
        if (!visited[v]) dfs1(v);
    order_arr[order_cnt++] = u;
}

// Kosaraju's Algorithm -- Pass 2: reverse DFS, assign component IDs
void dfs2(int u, int c) {
    comp[u] = c;
    for (int v : radj[u])
        if (comp[v] == -1) dfs2(v, c);
}

int main() {
    scanf("%d", &n);
    for (int i = 1; i <= n; i++) {
        int v;
        while (scanf("%d", &v) == 1 && v != 0) {
            adj[i].push_back(v);
            radj[v].push_back(i);
        }
    }

    // Pass 1
    memset(visited, false, sizeof(visited));
    order_cnt = 0;
    for (int i = 1; i <= n; i++)
        if (!visited[i]) dfs1(i);

    // Pass 2
    memset(comp, -1, sizeof(comp));
    int numSCC = 0;
    for (int i = order_cnt - 1; i >= 0; i--)
        if (comp[order_arr[i]] == -1)
            dfs2(order_arr[i], numSCC++);

    if (numSCC == 1) {
        printf("1\n0\n");
        return 0;
    }

    // Build condensed DAG: determine which SCCs have incoming/outgoing edges
    // We only need to know whether in-degree and out-degree are zero,
    // so duplicate edges between the same pair of SCCs do not affect the result.
    bool hasIn[MAXN] = {}, hasOut[MAXN] = {};
    for (int u = 1; u <= n; u++)
        for (int v : adj[u])
            if (comp[u] != comp[v]) {
                hasIn[comp[v]] = true;
                hasOut[comp[u]] = true;
            }

    int sources = 0, sinks = 0;
    for (int i = 0; i < numSCC; i++) {
        if (!hasIn[i]) sources++;
        if (!hasOut[i]) sinks++;
    }

    printf("%d\n%d\n", sources, max(sources, sinks));
    return 0;
}
\end{lstlisting}

\section{Complexity Analysis}

\begin{itemize}
  \item \textbf{Time complexity:} $O(n + m)$ where $m$ is the number of edges.
        Kosaraju's algorithm runs two DFS passes, each $O(n + m)$. Computing
        sources and sinks in the condensed DAG is $O(n + m)$.
  \item \textbf{Space complexity:} $O(n + m)$ for the adjacency lists and reverse graph.
\end{itemize}

\end{document}
